\begin{theorem}[Umbral singular en el índice $m$ (enunciado de cierre)]\label{thm:umbral_singular_indice_m_enunciado}
	Sea $\langle\cdot,\cdot\rangle_S$ un producto de Sobolev con componente continua soportada en
	$\overline{\mathbb D_R(0)}$, y sea $m=|I_{\mathrm{out}}|$ el número de átomos no-BPE.
	Suponemos además que estamos en régimen puramente no-BPE en la parte discreta, es decir,
	que todos los átomos discretos del producto corresponden a índices de $I_{\mathrm{out}}$.
	Supongamos que para todo $k\in I_{\mathrm{out}}$ el funcional
	\[
		\psi_k(p)=(\mathcal Dp)'(c_k)
	\]
	es no acotado respecto de $\|\cdot\|_S$. Entonces se verifica el patrón umbral
	\[
		\lim_{n\to\infty}\sigma_{j,n}=+\infty,
		\qquad 0\le j\le m-1,
	\]
	y
	\[
		\limsup_{n\to\infty}\sigma_{m,n}\le R.
	\]
	En particular, la inestabilidad queda confinada exactamente a $m$ direcciones singulares.
\end{theorem}

\begin{proof}
	Si $m=0$, la afirmación explosiva es vacía y la cota en el índice $m=0$ coincide con la
	estimación del caso puramente continuo. En adelante suponemos $m\ge 1$.

	\medskip
	\noindent
	\textbf{Paso 1: explosión de $\sigma_{j,n}$ para $0\le j\le m-1$.}

	Basta probarla para $j=m-1$, pues
	\[
		\sigma_{0,n}\ge\sigma_{1,n}\ge\cdots\ge\sigma_{m-1,n}
	\]
	para todo $n$.

	Fijemos $L>0$. Para cada $j\in I_{\mathrm{out}}$, por interpolación de Hermite en los
	puntos discretos $\{c_k\}_{k\in I_{\mathrm{out}}}$ existe $g_j\in\mathbb P[z]$ tal que
	\[
		g_j(c_j)=1,\quad g_j'(c_j)=0,
	\]
	y para $\ell\in I_{\mathrm{out}}$, $\ell\ne j$,
	\[
		g_j(c_\ell)=0,\quad g_j'(c_\ell)=0.
	\]

	Como la multiplicación por un polinomio fijo es continua en
	$(\mathbb P[z],\|\cdot\|_S)$, existe $C_j<\infty$ tal que
	\[
		\|g_jp\|_S\le C_j\|p\|_S,
		\qquad p\in\mathbb P[z].
	\]
	Como $\psi_j$ es no acotado, el Lema~\ref{lem:normalizacion_funcional_no_acotado}
	permite elegir $x_j\in\mathbb P[z]$ con
	\[
		\psi_j(x_j)=1,
		\qquad
		\|x_j\|_S<\frac{1}{L\,C_j\,\sqrt m}.
	\]
	Definimos $u_j:=g_jx_j$. Entonces
	\[
		\|u_j\|_S=\|g_jx_j\|_S\le C_j\|x_j\|_S<\frac{1}{L\sqrt m}.
	\]

	Además, por la regla del producto y las condiciones de Hermite,
	\[
		\psi_\ell(u_j)=\delta_{\ell j},
		\qquad \ell,j\in I_{\mathrm{out}}.
	\]
	En particular, los $u_j$ son linealmente independientes. Sea
	\[
		F:=\operatorname{span}\{u_j:\ j\in I_{\mathrm{out}}\},
	\]
	de dimensión $m$.

	Tomemos $p=\sum_{j\in I_{\mathrm{out}}}\alpha_ju_j\in F$. De la diagonalidad anterior,
	\[
		\psi_\ell(p)=\alpha_\ell,
		\qquad
		\Bigl(\sum_{\ell\in I_{\mathrm{out}}}|\psi_\ell(p)|^2\Bigr)^{1/2}=\|\alpha\|_{\ell^2}.
	\]
	Como los términos $|\psi_\ell(p)|^2=|(\mathcal Dp)'(c_\ell)|^2$ forman parte de la parte
	discreta de $\|\mathcal Dp\|_S^2$, se obtiene
	\[
		\|\mathcal Dp\|_S\ge \|\alpha\|_{\ell^2}.
	\]
	Por otro lado,
	\[
		\|p\|_S
		\le
		\sum_{j\in I_{\mathrm{out}}}|\alpha_j|\,\|u_j\|_S
		\le
		\|\alpha\|_{\ell^2}\Bigl(\sum_{j\in I_{\mathrm{out}}}\|u_j\|_S^2\Bigr)^{1/2}
		<
		\frac1L\|\alpha\|_{\ell^2}.
	\]
	Luego, para todo $p\in F\setminus\{0\}$,
	\[
		\frac{\|\mathcal Dp\|_S}{\|p\|_S}>L.
	\]

	Sea ahora
	\[
		n_L:=1+\max_{j\in I_{\mathrm{out}}}\deg(u_j).
	\]
	Si $n\ge n_L$, entonces $F\subset\mathbb P_{n-1}$ y por tanto
	$\mathcal Dp\in\mathbb P_n$ para todo $p\in F$; en consecuencia,
	$\mathbf D_n p=\mathcal Dp$ sobre $F$.

	Usando la forma max--min de Courant--Fischer para valores singulares (índices desde $0$),
	\[
		\sigma_{m-1,n}
		=
		\max_{\dim E=m}\ \min_{0\ne p\in E}\frac{\|\mathbf D_np\|_S}{\|p\|_S}
		\ge
		\min_{0\ne p\in F}\frac{\|\mathbf D_np\|_S}{\|p\|_S}
		=
		\min_{0\ne p\in F}\frac{\|\mathcal Dp\|_S}{\|p\|_S}
		> L.
	\]
	Como $L>0$ es arbitrario,
	\[
		\lim_{n\to\infty}\sigma_{m-1,n}=+\infty.
	\]
	Por monotonicidad en el índice, esto implica
	\[
		\lim_{n\to\infty}\sigma_{j,n}=+\infty,
		\qquad 0\le j\le m-1.
	\]

	\medskip
	\noindent
	\textbf{Paso 2: cota estable en el índice $m$.}

	Consideremos el subespacio canónico
	\[
		V_{\mathrm{out}}:=\bigcap_{k\in I_{\mathrm{out}}}\ker(\psi_k).
	\]
	Por independencia lineal de los funcionales $\{\psi_k\}_{k\in I_{\mathrm{out}}}$
	(interpolación de Hermite),
	\[
		\operatorname{codim}_{alg}(V_{\mathrm{out}})=m.
	\]
	Si $p\in V_{\mathrm{out}}$, entonces $(\mathcal Dp)'(c_k)=0$ para todo átomo discreto
	(activo) y, bajo la hipótesis puramente no-BPE, toda la parte discreta de
	$\|\mathcal Dp\|_S^2$ se anula. Por tanto,
	\[
		\|\mathcal Dp\|_S^2
		=
		\int_{\operatorname{supp}(\mu)}|z|^2|p(z)|^2\,d\mu(z)
		\le
		R^2\int |p(z)|^2\,d\mu(z)
		\le
		R^2\|p\|_S^2.
	\]
	Así,
	\[
		\|\mathcal D|_{V_{\mathrm{out}}}\|\le R.
	\]

	Definamos $V_{\mathrm{out},n}:=V_{\mathrm{out}}\cap\mathbb P_n$. Entonces
	\[
		\dim(V_{\mathrm{out},n})\ge n+1-m.
	\]
	Sea $\mathcal V_{\mathrm{out},n}\subset\mathbb C^{n+1}$ su imagen en coordenadas.
	Para $x\in\mathcal V_{\mathrm{out},n}$ asociado a $p\in V_{\mathrm{out},n}$,
	\[
		\|\mathbf D_nx\|_2
		=
		\|P_n(\mathcal Dp)\|_S
		\le
		\|\mathcal Dp\|_S
		\le
		R\|p\|_S
		=
		R\|x\|_2.
	\]
	Aplicando Courant--Fischer en su forma min--max,
	\[
		\sigma_{m,n}
		=
		\min_{\dim S=n+1-m}\ \max_{0\ne x\in S}\frac{\|\mathbf D_nx\|_2}{\|x\|_2}
		\le
		\max_{0\ne x\in\mathcal V_{\mathrm{out},n}}\frac{\|\mathbf D_nx\|_2}{\|x\|_2}
		\le R.
	\]
	En consecuencia,
	\[
		\limsup_{n\to\infty}\sigma_{m,n}\le R.
	\]
	Esto completa la prueba.
\end{proof}